\section{Introduction}


Modern object detectors employ many hand-crafted components~\citep{liu2020deep}, e.g., anchor generation, rule-based training target assignment, non-maximum suppression (NMS) post-processing. They are not fully end-to-end. Recently, \citet{carion2020end} proposed DETR to eliminate the need for such hand-crafted components, and built the first fully end-to-end object detector, achieving very competitive performance. DETR utilizes a simple architecture, by combining convolutional neural networks (CNNs) and Transformer~\citep{vaswani2017attention} encoder-decoders. They exploit the versatile and powerful relation modeling capability of Transformers to replace the hand-crafted rules, under properly designed training signals. 


Despite its interesting design and good performance, DETR has its own issues: (1) It requires much longer training epochs to converge than the existing object detectors. For example, on the COCO~\citep{lin2014microsoft} benchmark, DETR needs 500 epochs to converge, which is around 10 to 20 times slower than Faster R-CNN~\citep{ren2015faster}. (2) DETR delivers relatively low performance at detecting small objects. Modern object detectors usually exploit multi-scale features, where small objects are detected from high-resolution feature maps. Meanwhile, high-resolution feature maps lead to unacceptable complexities for DETR. The above-mentioned issues can be mainly attributed to the deficit of Transformer components in processing image feature maps. At initialization, the attention modules cast nearly uniform attention weights to all the pixels in the feature maps. Long training epoches is necessary for the attention weights to be learned to focus on sparse meaningful locations.
On the other hand, the attention weights computation in Transformer encoder is of quadratic computation w.r.t. pixel numbers. Thus, it is of very high computational and memory complexities to process high-resolution feature maps.


In the image domain, deformable convolution~\citep{dai2017deformable} is of a powerful and efficient mechanism to attend to sparse spatial locations. It naturally avoids the above-mentioned issues. While it lacks the element relation modeling mechanism, which is the key for the success of DETR.


In this paper, we propose \emph{Deformable DETR}, which mitigates the slow convergence and high complexity issues of DETR. It combines the best of the sparse spatial sampling of deformable convolution, and the relation modeling capability of Transformers. We propose the \emph{deformable attention module}, which attends to a small set of sampling locations as a pre-filter for prominent key elements out of all the feature map pixels.
The module can be naturally extended to aggregating multi-scale features, without the help of FPN~\citep{lin2017feature}. In Deformable DETR , we utilize (multi-scale) deformable attention modules to replace the Transformer attention modules processing feature maps, as shown in Fig.~\ref{fig:attention_demo}. 



Deformable DETR opens up possibilities for us to exploit variants of end-to-end object detectors, thanks to its fast convergence, and computational and memory efficiency. We explore a simple and effective \textit{iterative bounding box refinement} mechanism to improve the detection performance. We also try a \textit{two-stage Deformable DETR}, where the region proposals are also generated by a vaiant of Deformable DETR, which are further fed into the decoder for iterative bounding box refinement.

Extensive experiments on the COCO~\citep{lin2014microsoft} benchmark demonstrate the effectiveness of our approach. Compared with DETR, Deformable DETR can achieve better performance (especially on small objects) with 10$\times$ less training epochs. 
The proposed variant of two-stage Deformable DETR can further improve the performance.
Code is released at \url{https://github.com/fundamentalvision/Deformable-DETR}.


\begin{figure}[t]
\begin{center}
  \includegraphics[width=0.85\linewidth]{fig/overview.pdf}
\end{center}
\vspace{-0.5em}
\caption{Illustration of the proposed Deformable DETR object detector.}
\vspace{-0.5em}
\label{fig:attention_demo}
\end{figure}
